\chapter{1994年上海卷（文）}
\section{填空题}

\begin{problem}
不等式 \(|x + 1| < 1 \) 的解是 \fillinblank{}.

\end{problem}

\begin{problem}
函数 \(y = \sin^{2} x \) 的最小正周期是 \fillinblank{}.

\end{problem}

\begin{problem}
计算: \(\displaystyle\lim_{x\to1}\frac{x^{2}-2x-3}{x+1}= \) \fillinblank{}.

\end{problem}

\begin{problem}
以点 \(C(-2,3) \) 为圆心且与 \(y\) 轴相切的圆的方程是 \fillinblank{}.

\end{problem}

\begin{problem}
方程 \(\log_{3}(x-1)=\log_{9}(x+5) \) 的解是 \fillinblank{}.

\end{problem}

\begin{problem}
函数 \(y=\sqrt{x^{2}-1}(x\leqslant-1) \) 的反函数是 \fillinblank{}.

\end{problem}

\begin{problem}
双曲线 \(\frac{y^{2}}{2}-x^{2}=1 \) 的两个焦点的坐标是 \fillinblank{}.

\end{problem}

\begin{problem}
\(\left(x^{3}-\frac{2}{x^{2}}\right)^{5} \) 的展开式中 \(x^{5} \) 的系数是 (结果用数值表示) \fillinblank{}.

\end{problem}

\begin{problem}
如果复数 \(\frac{2+b\i}{1-2\i} \)（其中 \(\i\) 为虚数单位，\(b\) 为实数）的实部和虚部相等，那么 \(b=\) \fillinblank{}.

\end{problem}

\begin{problem}
某商场计划经销某种商品，预测可得盈利 30000 元、26000 元、8000 元的概率分别为 0.3、0.6、0.1，则经销该商品的期望盈利为 \fillinblank{} 元.

\end{problem}

\section{单选题}

\begin{problem}
当 \(a > 1\) 时，函数 \(y = \log_a x\) 和 \(y = (1 - a)x\) 的图象只可能是（\quad）

\choices
    {}
    {}
    {}
    {}

\end{problem}

\begin{problem}
如果 0 < a < 1，那么下列不等式中正确的是（\quad）

\choices
    {\((1 - a)^{\frac{1}{3}} > (1 - a)^{\frac{1}{2}}\)}
    {\(\log_{(1 - a)}(1 + a) > 0\)}
    {\((1 - a)^{3} > (1 + a)^{2}\)}
    {\((1 - a)^{1 + a} > 1\)}

\end{problem}

\begin{problem}
已知直线 \(2x + y - 2 = 0\) 和 \(mx - y + 1 = 0\) 的夹角为 \(\frac{\pi}{4}\)，那么 \(m\) 的值为（\quad）

\choices
    {\(-\frac{1}{3}\) 或 \(-3\)}
    {\(\frac{1}{3}\) 或 \(3\)}
    {\(-\frac{1}{3}\) 或 \(3\)}
    {\(\frac{1}{3}\) 或 \(-3\)}

\end{problem}

\begin{problem}
已知 \(a, b\) 是异面直线，直线 \(c\) 平行于直线 \(a\)，那么 \(c\) 与 \(b\)（\quad）

\choices
    {一定是异面直线}
    {一定是相交直线}
    {不可能是平行直线}
    {不可能是相交直线}

\end{problem}

\begin{problem}
设 \(l\) 是全集，集合 \(P, Q\) 满足 \(P \subset Q\)，则下面的结论中错误的是（\quad）

\choices
    {\(P \cup Q = Q\)}
    {\(\overline{P} \cup Q = I\)}
    {\(P \cap \overline{Q} = \emptyset\)}
    {\(\overline{P} \cap \overline{Q} = \overline{P}\)}

\end{problem}

\begin{problem}
已知 \(y = f(x) \left(x \in \R\right)\) 是奇函数，\(f(5) \neq 0\)，则下列各点中，在 \(y = f(x)\) 的图象上的点是（\quad）

\choices
    {\((5, f(-5))\)}
    {\((-5, -f(5))\)}
    {\((-5, f(5))\)}
    {\((5, -f(5))\)}

\end{problem}

\begin{problem}
设 \(a\)、\(b\) 是平面 \(a\) 外的任意两条线段，则“\(a\)、\(b\) 的长相等”是“\(a\)、\(b\) 在平面 \(a\) 内的射影长相等”的（\quad）

\choices
    {非充分条件也非必要条件}
    {充分必要条件}
    {必要条件而非充分条件}
    {充分条件而非必要条件}

\end{problem}

\begin{problem}
9 支足球队中，有 5 支亚洲队，4 支非洲队，从中任意抽取两队进行比赛，则 1 队是亚洲队且 1 队是非洲队的概率是（\quad）

\choices
    {\(\frac{\mathrm C_5^3 + \mathrm C_4^1}{\mathrm C_9^2}\)}
    {\(\frac{\mathrm C_4^1}{\mathrm C_9^2}\)}
    {\(\frac{\mathrm C_5^1}{\mathrm C_9^2}\)}
    {\(\frac{\mathrm C_5^1 \cdot \mathrm C_4^1}{\mathrm C_9^2}\)}

\end{problem}

\begin{problem}
在直角坐标系 \(xOy\) 中，曲线 \(C\) 的方程是 \(y = \cos x\)，现平移坐标系，把原点移到点 \(O' \left(\frac{\pi}{2}, -\frac{\pi}{2} \right)\)，则在坐标系 \(x'O'y'\) 中，曲线 \(C\) 的方程是（\quad）

\choices
    {\(y' = \sin x' + \frac{\pi}{2} \)}
    {\(y' = -\sin x' + \frac{\pi}{2} \)}
    {\(y' = \sin x' - \frac{\pi}{2} \)}
    {\(y' = -\sin x' - \frac{\pi}{2} \)}

\end{problem}

\begin{problem}
某个命题与自然数 \(n\) 有关. 如果当 \(n = k\)（\(k \in N\)）时该命题成立，那么可推得当 \(n = k + 1\) 时该命题也成立. 现已知当 \(n = 5\) 时该命题不成立，那么可推得（\quad）

\choices
    {当 \(n=6\) 时该命题不成立}
    {当 \(n=6\) 时该命题成立}
    {当 \(n=4\) 时该命题不成立}
    {当 \(n=4\) 时该命题成立}

\end{problem}

\section{解答题}

\begin{problem}
已知 \(\sin\alpha=\frac{3}{5},\alpha\in\left(\frac{\pi}{2},\pi\right),\mathrm{tg}(\pi-\beta)=\frac{1}{2} \) ，求 \(\mathrm{tg}(\alpha-2\beta) \) 的值.

\end{problem}

\begin{problem}
已知空间三个点 \(P(-2,0,2), Q(-1,1,2)\) 和 \(R(-3,0,4)\). 设 \(\bs{a} = \overrightarrow{PQ}, \bs{b} = \overrightarrow{PR}\). （1）求 \(\bs{a} \) 与 \(\bs{b} \) 的夹角 \(\theta \)（用反三角函数表示）；（2）试确定实数 \(k\)，使 \(k \bs{a} + \bs{b}\) 与 \(k \bs{a} - 2\bs{b}\) 互相垂直.

\end{problem}

\begin{problem}
如图，在梯形 \(ABCD\) 中，\(AD // BC\)，\(\angle ABC = \frac{\pi}{2}\)，\(AB = a\)，\(AD = 3a\)，且 \(\angle ADC = \arcsin \frac{\sqrt{3}}{5}\)，又 \(PA \perp\) 平面 \(ABCD\)，\(PA = a\). 求:

（1）二面角 \(P - CD - A\) 的大小（用反三角函数表示）；（2）点 \(A\) 到平面 \(PBC\) 的距离.

\end{problem}

\begin{problem}
如图，抛物线 \(y = 3x^{2}\) 和 \(y = -x^{2} + 4ax\ (a > 0)\) 交于原点 \(O\) 和点 \(A\).

（1）求这两条抛物线所围成的图形的面积；（2）过点 \(A\) 作抛物线 \(y = -x^{2} + 4ax\) 的切线，该切线与抛物线 \(y = 3x^{2}\) 交于另一点 \(P\)，求线段 \(AP\) 的垂直平分线的方程.

\end{problem}

\begin{problem}
设 \(\{a_{n}\}\) 是公差为 \(d\) 的等差数列，\(a_{3} + a_{5} = 2\)，\(S_{20} = a_{1} + a_{2} + \cdots + a_{20} = 150\)，又 \(b_{n} = 2^{a_{n} - 2a_{n+1}}\left(n \in N\right)\). （1）求 \(a_{1}, d\) 的值；（2）求证 \(\left|b_{n}\right|\) 是等比数列，并求 \(b_{4}\)；（3）设 \(k\) 为某个自然数，且满足 \(\displaystyle\lim_{n\to\infty}\left(b_k b_{k+1} + b_{k+1} b_{k+2} + \cdots + b_n b_{n+1}\right) = \frac{1}{96}\)，求 \(k\) 的值.

\end{problem}

\begin{problem}
设 \(P(1, t)\) 是直角坐标系内的点（\(t > 0\)）. 将向量 \(\overrightarrow{OP}\) 绕原点 \(O\) 按逆时针方向旋转 \(\frac{\pi}{2}\)，再把它的模变为原来的两倍，得到向量 \(\overrightarrow{OQ}\)，以 \(OP\)、\(OQ\) 为邻边作矩形 \(OPRQ\). （1）求点 \(Q\)、\(R\) 的坐标；

（2）求矩形 \(OPRQ\) 在第一象限部分的面积 \(S(t)\)；（3）确定函数 \(S(t)\) 的单调区间，并加以证明.
\end{problem}

